%% manuscript.tex
%% Improving four-class Fermi-LAT source classification through
%% feature engineering and covariate shift quantification
%%
%% Target: Monthly Notices of the Royal Astronomical Society (MNRAS)
%% Compile: pdflatex manuscript (twice)

\documentclass[fleqn,12pt]{article}

\usepackage{graphicx}
\usepackage{amsmath,amssymb,bm}
\usepackage{natbib}
\usepackage{url}
\usepackage{booktabs}
\usepackage{longtable}
\usepackage{multirow}
\usepackage{xcolor}

\usepackage[margin=2.5cm]{geometry}
\usepackage{setspace}
\doublespacing
\usepackage{lineno}
\linenumbers

% ── MNRAS-style shortcuts ─────────────────────────────────────────────
\newcommand{\fermi}{\textit{Fermi}-LAT}
\newcommand{\fgl}{4FGL-DR4}
\newcommand{\Gam}{\Gamma}
\newcommand{\sigs}{\sigma_{\rm curv}}
\newcommand{\GamLP}{\Gamma_{\rm 1\,GeV}}
\newcommand{\betaLP}{\beta_{\rm LP}}
\newcommand{\Epeak}{E^{\rm LP}_{\rm peak}}
\newcommand{\logV}{\log V}
\newcommand{\GALL}{|b|}
\newcommand{\Fvar}{F_{\rm var}}
\newcommand{\ARI}{\mathrm{ARI}}
\newcommand{\BAcc}{\mathrm{BAcc}}
\newcommand{\malyshev}{Malyshev \& Bhat}

%% ────────────────────────────────────────────────────────────────────
\title{Improving four-class \textit{Fermi}-LAT source classification through
 spectral feature engineering and covariate shift quantification}

\author{Debarshi Mukherjee$^{1}$\thanks{E-mail: debarshi99mukherjee@gmail.com}
\and Pratik Majumdar$^{2}$\\[4pt]
$^{1}$Department of Physical Sciences, IISER Berhampur, Berhampur 760003, India\\
$^{2}$Saha Institute of Nuclear Physics (HBNI), Kolkata 700064, India}

\date{\today}


\begin{document}
\maketitle

%% ── Abstract ─────────────────────────────────────────────────────────
\begin{abstract}
We present a machine learning classification study of the 2,423
unassociated gamma-ray sources in the \fermi{} 14-year Fourth Source
Catalogue (\fgl{}; Ballet et al.\ 2023), discriminating between four
source classes: BL~Lacertae objects (BLL), Flat Spectrum Radio Quasars
(FSRQ), young pulsars (PSR), and millisecond pulsars (MSP). Building on
the seven-feature baseline of \malyshev{} (2023), we introduce three
modifications that together raise the 10-fold cross-validated balanced
accuracy of a Random Forest classifier from $84.5 \pm 3.3\,\%$ to
$86.5 \pm 3.4\,\%$ on the same labeled sample of 2,630 sources
(Wilcoxon signed-rank test: $p = 0.029$). The core modification is the
replacement of the naive power-law index \texttt{PL\_Index} with the
log-parabola spectral index evaluated at 1\,GeV,
$\GamLP \equiv \Gamma_{\rm LP} + 2\betaLP\,\ln(1\,\text{GeV}/E_{\rm piv})$
\citep{Malyshev2023}, supplemented by the log-parabola peak energy
$\log\Epeak$ and the fractional variability $\Fvar$. The gain is
concentrated in the FSRQ class ($\Delta F_1 = +0.026$), whose
separation from BLL improves when spectral curvature is measured at a
physically consistent reference energy rather than at the individually
fitted pivot energy; a direct $\GamLP$--$\log\Epeak$ projection confirms
that this gain traces the classical blazar sequence. We stress-test
the feature-set comparison with repeated cross-validation (5 random
seeds $\times$ 10 folds, 50 paired estimates) and independent 70/30
hold-out resampling (15 repeats), both of which confirm the gain at
higher significance ($p = 8.8\times10^{-4}$ and $p = 0.022$,
respectively) than the single-split test alone. Permutation importance
corroborates the Gini-based feature ranking, with the log-parabola
curvature $\betaLP$ and the spectral index $\GamLP$ remaining the two
most informative features under both metrics. We quantify the
covariate shift between labeled and unassociated sources across all
ten features using Wasserstein distances and KS tests, and find that a
binary shift classifier achieves AUC\,$= 0.985$, confirming that the
two populations are nearly perfectly separable in feature space. Of
the 2,423 unassociated sources, 34.0\% fall below the 5th percentile
of the training significance distribution and are designated
out-of-distribution (OOD); their classifications carry reduced
reliability and are clearly flagged in our published table. An
out-of-fold calibration check shows that our adopted high-confidence
threshold, $P_{\rm max} \geq 0.80$, corresponds to an observed accuracy
of 96.1\,\% among in-distribution sources, validating its use for
follow-up prioritisation. Applying the 10-feature Random Forest
(retrained deterministically; see Data Availability) to the
1,598 remaining in-distribution sources, we identify 541 high-confidence
candidates ($P_{\rm max} \geq 0.80$): 274 BLL, 61 FSRQ, 74 PSR, and
132 MSP. All posterior probability tables, the trained model, and the
analysis and validation code are provided as supplementary
machine-readable data.
\end{abstract}

\noindent\textbf{Key words:}~
gamma-ray astronomy --- catalogues --- methods: statistical ---
BL Lacertae objects: general --- pulsars: general ---
quasars: general --- surveys
\\[6pt]

%% ── 1. Introduction ──────────────────────────────────────────────────
\section{Introduction}
\label{sec:intro}

The \fermi{} Large Area Telescope (LAT) has operated continuously since
2008 \citep{Atwood2009}, surveying the entire sky above 50\,MeV. Its
source catalogues have been updated through a succession of releases,
from the Second \fermi{} LAT Catalog \citep{Nolan2012} through the
Third Catalog \citep{Acero2015_3FGL}, the Fourth Catalog
\citep{Abdollahi2020} and its associated catalogue of
active galactic nuclei \citep{Ackermann2015_3LAC,Ajello2020}. The most
recent source list, \fgl{} \citep{Ballet2023}, contains 7,195
gamma-ray sources. Among
these, roughly 30 per cent -- 2,423 sources -- carry no association
with any catalogued counterpart at other wavelengths.
Identifying the likely source classes among these unassociated
detections advances both population statistics and the scheduling of
radio and X-ray follow-up campaigns, such as the Swift-XRT/UVOT
multiwavelength surveys that have independently characterised
counterparts to hundreds of unassociated sources \citep{Kerby2021}.

The dominant identified classes are extragalactic: 1,490 BL~Lacertae
objects (BLL), 820 Flat Spectrum Radio Quasars (FSRQ), and 1,623
blazar candidates of uncertain type (BCU). The second-most important
identified population is Galactic: 141 young pulsars (PSR) and 179
millisecond pulsars (MSP), as catalogued in the Third \fermi{} Pulsar
Catalog \citep{Smith2023}, building on the earlier Second \fermi{}
Pulsar Catalog \citep{Abdo2013} and the substantial MSP discoveries
enabled by radio timing campaigns targeting \fermi{} unassociated
sources \citep{Guillemot2016}. Because the four confirmed classes -- BLL,
FSRQ, PSR, MSP -- differ profoundly in their underlying physics and
their multi-wavelength counterpart types, a reliable four-class
classification is considerably more informative for follow-up work
than the binary AGN/non-AGN separation that has dominated the
literature \citep{Ackermann2012,Mirabal2012,Doert2014}.

Machine learning classification of \fermi{}-LAT sources has been
studied extensively. \citet{Chiaro2016} applied artificial neural
networks to separate BLL from FSRQ in 3FGL, an approach later refined
by \citet{Lefaucheur2017} using multivariate boosted-tree and
neural-network classifiers on 3FGL blazar candidates. \citet{Kang2019} and
\citet{Kovacevic2020} obtained accuracies of $91$--$93\,\%$ for
binary blazar classification on 4FGL. \citet{Luo2020} systematically
investigated how feature selection and sample size affect
classification performance, motivating the compact feature sets used
in later work, while \citet{Finke2021} explored classification
directly from unbinned energy and time spectra with deep learning.
\citet{Zhu2021} and \citet{CoronadoBlazquez2022} extended
multi-class and gradient-boosted classification to the full
4FGL/4FGL-DR3 unassociated population, and \citet{Bhat2022}
constructed probabilistic multi-method catalogues combining Random
Forest, boosted trees, and neural-network posteriors. \citet{Cooper2023} reached
$91.2\,\%$ using SuperLearner ensembles on 4FGL-DR3, and
\citet{Zhu2024} used a multi-method voting strategy on the same
release, demonstrating that ensemble approaches improve stability.
\citet{Bhatta2024} explored self-supervised pre-training on 4LAC-DR3.
Three-class and four-class problems -- adding pulsars alongside the
two blazar types -- were addressed by \citet{Saz_Parkinson2016} on
3FGL ($\sim$70\,\% balanced accuracy) and by \citet{Zhu2024} on
4FGL-DR3 ($\sim$72\,\%).

Two methodological issues have received increasing attention but
remain incompletely resolved. First, the choice of spectral index
feature matters more than is commonly appreciated.
\citet{Malyshev2023} noted that the naive power-law index
\texttt{PL\_Index}, computed for all sources regardless of their
preferred spectral model, is a suboptimal discriminator because for
sources best described by a log-parabola (LP) or sub-exponential
cutoff (PLEC) spectrum, \texttt{PL\_Index} conflates two physically
distinct quantities: the intrinsic low-energy slope and the mean
slope artificially flattened by spectral curvature. Their prescription
of evaluating the LP spectral index at a common reference energy of
1\,GeV, $\GamLP$, partially resolves this issue for blazar-like
sources. We show here that supplementing $\GamLP$ with the peak
energy of the LP SED, $\log\Epeak$, provides additional separation
between BLL and FSRQ.

Second, the labeled training set (associated sources) and the
classification target (unassociated sources) are not drawn from the
same distribution. \citet{MalyshevRASTI2023} quantified this covariate
shift and its effect on classification performance for the 4FGL
sample; we extend this analysis with an explicit per-feature
Wasserstein/KS characterisation and a binary shift-classifier AUC for
\fgl{}, following the framework of \citet{Shimodaira2000}.

The present paper makes the following specific contributions:
(i) a systematic, statistically tested comparison of the
\citet{Malyshev2023} seven-feature set with a ten-feature extension
on the full \fgl{} labeled sample;
(ii) a complete, per-feature covariate shift characterisation for
\fgl{};
(iii) an explicit OOD designation for unassociated sources that fall
outside the support of the training distribution;
and (iv) a probabilistic four-class classification table for all 2,423
unassociated \fgl{} sources, with per-source confidence scores and OOD
flags, made available as a machine-readable catalogue.

We do not claim to supersede \citet{Malyshev2023} or \citet{Zhu2024};
rather, we extend both studies to \fgl{} with a rigorously
tested feature set and a transparent handling of distribution shift.

The paper is organised as follows. Section~\ref{sec:data} describes the
data and feature construction. Section~\ref{sec:shift} presents the
covariate shift analysis. Section~\ref{sec:method} describes the
classifiers and validation strategy. Results are given in
Section~\ref{sec:results}. Section~\ref{sec:disc} discusses the
physical interpretation and limitations. Section~\ref{sec:conc}
summarises our conclusions.

%% ── 2. Data and Feature Construction ─────────────────────────────────
\section{Data and feature construction}
\label{sec:data}

\subsection{The 4FGL-DR4 catalogue}
\label{ssec:4fgl}

We use the \fgl{} catalogue \citep{Ballet2023} in its FITS release
\texttt{gll\_psc\_v35.fit}, downloaded from the Fermi Science Support
Center. The main source table contains 7,195 rows and 79 spectral,
positional, variability, and association columns. We read this file
with a bespoke pure-Python FITS binary-table parser (Appendix~A) in
environments where \textsc{astropy} is unavailable, verified against
the published source statistics.

The labeled sample comprises all sources with a non-blank
\texttt{CLASS1} field whose stripped lower-case value falls in
$\{$\texttt{bll, BLL, fsrq, FSRQ, psr, PSR, msp, MSP}$\}$: a total
of 2,630 sources (BLL: 1,490; FSRQ: 820; PSR: 141; MSP: 179). Both
lower-case (likely association) and upper-case (firm identification)
entries are included, since the association probability for
lower-case classes in \fgl{} exceeds 80 per cent by construction
\citep{Ballet2023}. Blazar candidates of uncertain type
(\texttt{bcu}/\texttt{BCU}; 1,623 sources) are treated as a
separate prediction-only set. The 2,423 sources with a blank
\texttt{CLASS1} field are the primary classification target.

\subsection{Feature set}
\label{ssec:features}

Our ten features are listed in Table~\ref{tab:features} and motivated
below.

The first seven features follow \citet{Malyshev2023} exactly:
$\log F_E$ (log$_{10}$ energy flux in erg\,cm$^{-2}$\,s$^{-1}$),
$\log\sigma_{F_E}$ (its log uncertainty),
$\log\sqrt{TS}$ (log detection significance),
$\GamLP$ (LP spectral index at 1\,GeV; see below),
$\betaLP$ (log-parabola curvature parameter),
$\sigs$ (log-parabola curvature significance),
and $\logV$ (log variability index).
Three features are added by the present work:
$\GALL$ (absolute Galactic latitude),
$\log\Epeak$ (log peak energy of the LP SED),
and $\Fvar$ (fractional variability amplitude).

\paragraph{The $\GamLP$ feature.}
The \fgl{} catalogue fits three spectral models to each source: a
power law (PL), a log-parabola (LP), and a sub-exponential cutoff
power law (PLEC). The column \texttt{PL\_Index} contains the PL
best-fit index regardless of which model is preferred. For sources
whose preferred model is LP or PLEC, this column conflates the
intrinsic low-energy slope with the effect of spectral curvature. For
the LP model,
\begin{equation}
  \frac{dN}{dE} \propto
  \left(\frac{E}{E_{\rm piv}}\right)^{-\Gamma_{\rm LP}
    - \betaLP\ln(E/E_{\rm piv})} ,
  \label{eq:lp}
\end{equation}
the local spectral index at energy $E$ is
$\Gamma(E) = \Gamma_{\rm LP} + 2\betaLP\,\ln(E/E_{\rm piv})$.
Evaluating this at a common reference energy $E_0 = 1\,\text{GeV}$
gives
\begin{equation}
  \GamLP \equiv \Gamma_{\rm LP}
    + 2\betaLP\,\ln\!\left(\frac{1\,\text{GeV}}{E_{\rm piv}}\right) ,
  \label{eq:gamma1gev}
\end{equation}
which is directly comparable across sources with different pivot
energies. This definition was introduced by \citet{Malyshev2023} and
we adopt it unchanged.

\paragraph{The $\log\Epeak$ feature.}
For a log-parabola spectrum the SED peak in $\nu F_\nu$ occurs at
$\Epeak = E_{\rm piv}\exp[-(2-\Gamma_{\rm LP})/(2\betaLP)]$,
available in \fgl{} as column \texttt{LP\_EPeak}. This quantity is
physically distinct from $\GamLP$: BLL with high-synchrotron-peak
emission have larger $\Epeak$ than low-synchrotron-peak BLL and FSRQ,
providing a separation that complements the spectral index
(see Fig.~\ref{fig:param_space}, panel b). For the 292 sources
(11.1\,\%) lacking a defined \texttt{LP\_EPeak} (those fit by a pure
power law with negligible curvature), we impute the column median
(738.6\,MeV); this fraction is comparable across all four classes
and therefore does not introduce class-specific information.

\paragraph{The $\Fvar$ feature.}
The fractional variability amplitude $\Fvar$, normalised to the mean
flux level, provides information orthogonal to the variability index
$V$: $V$ measures the total chi-squared-like variability significance
while $\Fvar$ reflects the relative amplitude, distinguishing between
a source that is significantly variable but only mildly so in
amplitude from one that is both significant and large-amplitude. FSRQs
exhibit higher $\Fvar$ than BLL on average, consistent with their
stronger disk--jet coupling \citep{Abdo2010_3LAC}.

\paragraph{The Galactic latitude feature.}
Young pulsars are strongly concentrated toward the Galactic plane
($|b| \lesssim 5^\circ$; \citealt{Abdo2013}), while millisecond
pulsars have a broader but still plane-preferring distribution
($|b| \lesssim 20^\circ$; \citealt{Guillemot2016}). Blazars are
distributed approximately isotropically on the sky. Absolute Galactic
latitude $\GALL$ is therefore a powerful PSR/blazar discriminator that
is not included in the baseline feature set of \citet{Malyshev2023}.

\begin{table}
\caption{The ten-feature set used in this work.
The first seven reproduce the baseline of \protect\citet{Malyshev2023}.
$^a$~Definition given in equation~(\ref{eq:gamma1gev}).
$^b$~Median-imputed for the 11.1\,\% of sources lacking a defined
\texttt{LP\_EPeak}.}
\label{tab:features}
\begin{tabular}{lll}
\hline
Symbol & \fgl{} column & Physical content \\
\hline
$\log F_E$       & \texttt{log(Energy\_Flux100)} & Energy flux at 100 MeV \\
$\log\sigma_{F}$ & \texttt{log(Unc\_Energy\_Flux100)} & Flux uncertainty \\
$\log\sqrt{TS}$  & \texttt{log(Signif\_Avg)}  & Detection significance \\
$\GamLP$$^a$     & \texttt{LP\_Index, LP\_beta,} & LP index at 1\,GeV \\
                 & \texttt{Pivot\_Energy}       &  \\
$\betaLP$        & \texttt{LP\_beta}            & LP curvature \\
$\sigs$          & \texttt{LP\_SigCurv}         & Curvature significance \\
$\logV$          & \texttt{log(Variability\_Index)} & Variability \\
$\GALL$          & \texttt{abs(GLAT)}           & Galactic latitude \\
$\log\Epeak$$^b$ & \texttt{log(LP\_EPeak)}      & LP SED peak energy \\
$\Fvar$          & \texttt{Frac\_Variability}   & Fractional variability \\
\hline
\end{tabular}
\end{table}

Fig.~\ref{fig:param_space} illustrates the four-class structure
in four two-dimensional projections of the feature space. The
curvature significance $\sigs$ cleanly separates pulsars
(both PSR and MSP, which show high curvature due to their spectral
cutoff at 1--5\,GeV) from blazars (panel a); Galactic latitude
separates young pulsars from all other classes (panel c); and
$\GamLP$--$\betaLP$ separates FSRQ from BLL within the blazar
population (panel b).

\begin{figure}
\includegraphics[width=\columnwidth]{figs_final/fig1_parameter_space.pdf}
\caption{Four projections of the ten-dimensional feature space for
  the 2,630 labeled \fgl{} sources (BLL: blue; FSRQ: red; PSR: green;
  MSP: purple). (a)~$\GamLP$ versus $\sigs$: pulsars separate cleanly
  from blazars via spectral curvature. (b)~$\GamLP$ versus $\betaLP$:
  FSRQ occupy the high-$\betaLP$, soft-$\Gamma$ region. (c)~$\GALL$
  versus $\sigs$: PSR concentrate at low latitude. (d)~$\logV$ versus
  $\Fvar$: FSRQ display systematically larger variability amplitude
  than BLL.}
\label{fig:param_space}
\end{figure}

\begin{figure*}
\includegraphics[width=\textwidth]{figs_final/fig2_feature_distributions.pdf}
\caption{Marginal distributions of all ten features for each of the
  four labeled source classes. Dashed vertical lines mark the class
  median. $\sigs$ provides the strongest single-feature pulsar--blazar
  discrimination; $\GamLP$ provides the strongest single-feature
  BLL--FSRQ discrimination.}
\label{fig:distributions}
\end{figure*}

Table~\ref{tab:class_stats} summarises the median and
interquartile range of each feature by class.

\begin{table*}
\caption{Median (interquartile range) of each feature for the four
  labeled source classes in \fgl{}.}
\label{tab:class_stats}
\begin{tabular}{lcccc}
\hline
Feature & BLL ($N$\,=\,1490) & FSRQ ($N$\,=\,820)
        & PSR ($N$\,=\,141)  & MSP ($N$\,=\,179) \\
\hline
$\GamLP$       & $1.84\;(0.36)$ & $2.45\;(0.30)$ & $2.11\;(0.37)$ & $1.90\;(0.42)$ \\
$\betaLP$      & $0.09\;(0.10)$ & $0.18\;(0.13)$ & $0.15\;(0.13)$ & $0.12\;(0.11)$ \\
$\sigs$        & $1.49\;(2.45)$ & $2.37\;(4.21)$ & $11.51\;(14.5)$& $6.55\;(8.78)$ \\
$\logV$        & $1.49\;(0.75)$ & $2.15\;(0.98)$ & $1.17\;(0.51)$ & $1.18\;(0.51)$ \\
$\Fvar$        & $0.14\;(0.22)$ & $0.40\;(0.45)$ & $0.02\;(0.05)$ & $0.02\;(0.04)$ \\
$\log\Epeak$   & $3.17\;(0.73)$ & $2.08\;(0.46)$ & $2.73\;(0.61)$ & $2.90\;(0.65)$ \\
$\GALL$ (deg)  & $37.6\;(44.3)$ & $28.8\;(38.1)$ & $2.0\;(3.3)$   & $22.3\;(33.7)$ \\
$\log F_E$     & $-11.60\;(0.61)$&$-11.55\;(0.65)$&$-11.80\;(0.71)$&$-12.05\;(0.67)$\\
$\log\sqrt{TS}$& $1.26\;(0.46)$ & $1.30\;(0.44)$ & $1.48\;(0.51)$ & $1.09\;(0.37)$ \\
$\log\sigma_F$ & $-12.07\;(0.63)$&$-12.07\;(0.65)$&$-12.32\;(0.71)$&$-12.63\;(0.72)$\\
\hline
\end{tabular}
\end{table*}

%% ── 3. Covariate Shift ───────────────────────────────────────────────
\section{Covariate shift between labeled and unassociated sources}
\label{sec:shift}

Training a classifier on the labeled subsample and deploying it on
unassociated sources implicitly assumes that both sets are drawn from
the same marginal feature distribution $p(\mathbf{x})$. This
assumption is violated in flux-limited surveys: unassociated sources
are, by selection, those with either a faint or atypical counterpart,
or no catalogued counterpart at all. They are therefore systematically
fainter, lower-significance, and distributed differently on the sky
than the associated population.

We quantify this shift for each feature using two complementary
statistics. The \emph{Wasserstein $W_1$ distance} (Earth Mover's
distance) measures the minimum transport cost to morph one empirical
distribution into the other; it is sensitive to both location and
shape differences and carries physical units. The \emph{Kolmogorov--Smirnov
(KS) statistic} $D_{\rm KS}$ \citep{Massey1951} tests the null hypothesis that the two
samples share a common distribution. Statistical significance of
paired classifier comparisons throughout this paper is assessed with
the Wilcoxon signed-rank test \citep{Wilcoxon1945}.

The results are listed in Table~\ref{tab:shift}. All ten features show
highly significant shift ($p < 10^{-22}$ for all features), consistent
with \citet{MalyshevRASTI2023} who found significant covariate shift
and its effect on classification performance in \fgl{}. The largest
Wasserstein distances are for Galactic
latitude ($W_1 = 17.9$\,deg), reflecting the relative excess of
low-latitude sources among unassociated detections, and for the
variability index ($W_1 = 0.50$\,dex in log-scale). The largest
relative shift in KS terms is for $\log\sqrt{TS}$ ($D_{\rm KS} = 0.567$),
because the detection significance of unassociated sources is
systematically lower than for associated ones: the unassociated
median $\log\sqrt{TS} = 0.74$ versus a labeled median of $1.26$.

\begin{table}
\caption{Covariate shift between labeled ($N = 2630$) and unassociated
  ($N = 2423$) \fgl{} sources. $W_1$ is the Wasserstein $L^1$ distance
  in the units of each feature (after log-transformation where
  applicable). All $p$-values are far below any conventional threshold;
  we cite exact values to compare the \emph{magnitude} of shift across
  features.}
\label{tab:shift}
\begin{tabular}{lcc}
\hline
Feature & $W_1$ & $D_{\rm KS}$ \\
\hline
$\log F_E$         & 0.237 & 0.176 \\
$\log\sigma_F$     & 0.142 & 0.211 \\
$\log\sqrt{TS}$    & 0.442 & 0.567 \\
$\GamLP$           & 0.161 & 0.159 \\
$\betaLP$          & 0.091 & 0.341 \\
$\sigs$            & 1.685 & 0.141 \\
$\logV$            & 0.501 & 0.447 \\
$\GALL$ (deg)      & 17.89 & 0.400 \\
$\log\Epeak$       & 0.288 & 0.145 \\
$\Fvar$            & 0.263 & 0.385 \\
\hline
\multicolumn{3}{l}{Shift classifier AUC = 0.985}\\
\hline
\end{tabular}
\end{table}

To give an integrated measure of how separable the two populations are,
we train a Random Forest to distinguish labeled from unassociated
sources in the ten-dimensional feature space. The resulting AUC of
$0.985$ indicates near-perfect separability -- i.e.\ the joint feature
distribution is almost completely different between the populations.
This is a considerably more severe shift than is typical in machine
learning benchmarks, and it implies that naive application of a
classifier trained on the labeled sample will produce biased posterior
probabilities for the unassociated sources.

Fig.~\ref{fig:shift} illustrates the per-feature distributions,
showing the full shape difference for each of the ten features.

\begin{figure*}
\includegraphics[width=\textwidth]{figs_final/fig3_covariate_shift.pdf}
\caption{Marginal feature distributions for labeled (blue) and
  unassociated (orange) sources. The Wasserstein $W_1$ distance and
  KS statistic $D_{\rm KS}$ are annotated in each panel. All features
  show significant distributional differences; detection significance
  ($\log\sqrt{TS}$) and Galactic latitude ($\GALL$) show the most
  pronounced shift.}
\label{fig:shift}
\end{figure*}

Following \citet{Shimodaira2000}, we flag any unassociated source with
$\log\sqrt{TS}$ below the 5th percentile of the training distribution
as out-of-distribution (OOD). This threshold of $\log\sqrt{TS} = 0.69$
identifies 825 sources (34.0\,\% of the unassociated sample). For
these sources, the classifier is extrapolating to a significance
regime it has barely seen in training, and the resulting posterior
probabilities should be treated with caution. We retain their
predicted labels in the published table but clearly mark them.

%% ── 4. Classification Methods ────────────────────────────────────────
\section{Classification methods}
\label{sec:method}

\subsection{Classifiers}

We focus on the Random Forest \citep[RF;][]{Ho1995,Breiman2001} as our
production classifier,
since it (i) achieves the highest or joint-highest cross-validated
balanced accuracy among the seven classifiers we evaluate,
(ii) provides interpretable feature importances, and (iii) is
directly comparable to \citet{Malyshev2023} who also use RF as their
primary model.

The RF is configured with 500 trees, $\sqrt{p}$ features per split,
\texttt{min\_samples\_leaf}\,$= 2$, and balanced class weights to
counteract the PSR and MSP minority classes. All feature matrices
are standardised to zero mean and unit variance using a
\texttt{StandardScaler} fit on the training fold only to prevent
data leakage.

We also train and evaluate a Gradient Boosting classifier
\citep[GB;][]{Friedman2001}
(200 estimators, learning rate 0.08, max depth 4, subsample 0.8) and
a Multi-Layer Perceptron \citep[MLP;][]{Rumelhart1986}
(architecture $10\to128\to64\to32\to4$;
Adam optimiser \citep{Kingma2015}; early stopping; 12,708 trainable
parameters) for
comparison. The MLP architecture is illustrated in
Fig.~\ref{fig:mlp_arch}.

\begin{figure}
\includegraphics[width=\columnwidth]{figs_final/fig10_mlp_architecture.png}
\caption{Architecture of the Multi-Layer Perceptron classifier: a
  10-node standardised input layer (one node per feature in
  Table~\ref{tab:features}), three ReLU-activated hidden layers of
  128, 64, and 32 units, and a 4-node softmax output layer
  corresponding to the BLL/FSRQ/PSR/MSP classes. The network is
  trained with the Adam optimiser \citep{Kingma2015}, $\ell_2$
  regularisation $\alpha=10^{-4}$, a batch size of 64, and early
  stopping on a 15\,\% internal validation split (20 iterations'
  patience), giving 12,708 trainable parameters in total. For large
  hidden layers only a representative subset of nodes and connections
  is drawn, for legibility.}
\label{fig:mlp_arch}
\end{figure}

\subsection{Validation strategy}

Performance is assessed in two complementary ways. For the
feature-set comparison (Section~\ref{sec:results_feat}), we use
10-fold stratified cross-validation on the full labeled sample of
2,630 sources and report mean $\pm$ standard deviation of the
balanced accuracy. The balanced accuracy,
$\BAcc = \frac{1}{4}\sum_{k}\mathrm{Recall}_k$, is the appropriate
metric here because of the factor of $\sim$10 imbalance between
the dominant BLL class and the PSR and MSP minorities. For the
final classifier comparison (Section~\ref{sec:results_clf}), we use
a single stratified 70/30 hold-out split (1,839 train, 791 test).
All reported metrics other than the cross-validated comparison are
evaluated on this held-out test set.

Statistical significance of the feature-set gain is assessed with a
paired Wilcoxon signed-rank test on the fold-level balanced accuracy
scores of the 7-feature and 10-feature RF, ensuring identical random
seeds and data folds.

\subsection{Robustness checks}
\label{ssec:robustness}

A single 10-fold CV run or a single 70/30 split can be sensitive to
the particular random partition chosen, especially with the
moderate-sized minority classes (PSR: $N=141$; MSP: $N=179$) in our
sample. We therefore repeat the cross-validated feature-set comparison
over 5 independent random seeds (50 paired fold-level estimates in
total) and repeat the 70/30 hold-out evaluation over 15 independent
stratified splits, reporting the mean, standard deviation, and paired
Wilcoxon $p$-value across repeats in each case
(Section~\ref{sec:results_feat}).

We additionally compute permutation importance
\citep{Breiman2001} on the held-out test fold (15 repeats), which
measures the balanced-accuracy loss from randomly shuffling one
feature at a time. Unlike Gini importance, permutation importance is
insensitive to the number of ways a continuous feature can produce a
split, and comparing the two rankings guards against over-interpreting
a Gini-based ranking that is inflated by that effect
(Section~\ref{sec:results_fi}).

Finally, we assess the calibration of the RF's posterior probabilities
\citep{Platt1999,NiculescuMizil2005}
by pooling out-of-fold predictions across a 10-fold CV split and
comparing the predicted $P_{\rm max}$ in each probability bin against
the observed fraction of correct classifications. This directly tests
whether our adopted high-confidence threshold of $P_{\rm max}\geq0.80$
corresponds to the accuracy level its name implies
(Section~\ref{sec:results_unassoc}).

%% ── 5. Results ───────────────────────────────────────────────────────
\section{Results}
\label{sec:results}

\subsection{Effect of the extended feature set}
\label{sec:results_feat}

Fig.~\ref{fig:feat_compare} shows the main result: the 10-feature RF
achieves $\BAcc = 86.5 \pm 3.4\,\%$ in 10-fold CV, compared to
$84.5 \pm 3.3\,\%$ for the 7-feature baseline on the same 2,630
sources, a gain of $+2.3$ percentage points that is statistically
significant (Wilcoxon $p = 0.029$). The gain is primarily carried by
the FSRQ class ($\Delta F_1 = +0.026$) and, to a lesser extent, BLL
($\Delta F_1 = +0.006$). The PSR and MSP F1 scores change by less
than 0.005, indicating that the three new features improve blazar
sub-classification rather than pulsar identification.

\begin{figure}
\includegraphics[width=\columnwidth]{figs_final/fig4_feature_comparison.pdf}
\caption{Feature-set comparison. \emph{Left:} 10-fold CV balanced
  accuracy for three classifiers under the 7-feature \citet{Malyshev2023}
  baseline and the 10-feature set of this work. Error bars show
  $\pm 1\sigma$ over folds. The RF improvement is statistically
  significant (Wilcoxon $p = 0.029$). \emph{Right:} Per-class $F_1$
  gain (10-feature minus 7-feature RF) on the held-out test set. The
  gain concentrates in the FSRQ class.}
\label{fig:feat_compare}
\end{figure}

The physical interpretation is straightforward, and traces back to the
orientation-based unified schemes for radio-loud AGN
\citep{UrryPadovani1995}, in which BLL and FSRQ are understood as the
same underlying jetted population viewed at different accretion
states. The primary gain comes
from $\log\Epeak$: FSRQ tend to have lower-energy LP SED peaks than
BLL (Table~\ref{tab:class_stats}: FSRQ median $\log\Epeak = 2.08$
versus BLL $3.17$), consistent with the blazar sequence
\citep{Fossati1998,GhiselliniTavecchio2008,Ghisellini2017}, and this variable adds separation
orthogonal to $\GamLP$ and $\betaLP$ alone. Fig.~\ref{fig:blazarseq}
makes this connection explicit: plotting $\GamLP$ against $\log\Epeak$
for the labeled sample reveals the classical blazar-sequence trend --
sources with softer 1\,GeV index cluster at lower peak energy, and the
two features together separate FSRQ from BLL considerably more
cleanly than either alone, in a projection that is directly comparable
to the classic $\alpha_{\rm ox}$--luminosity sequence of
\citet{Fossati1998}. The $\Fvar$ feature adds a modest but consistent
improvement by distinguishing high-amplitude FSRQ variability from the
lower-amplitude BLL variability, even when both have high variability
index $V$. The $\GALL$ feature adds almost nothing for blazar
sub-classification but prevents a small number of Galactic-plane
sources from being misclassified as extragalactic.

\begin{figure}
\includegraphics[width=\columnwidth]{figs_final/fig9_blazar_sequence.pdf}
\caption{Blazar-sequence projection of the two features responsible
  for the largest share of the feature-set gain: the 1\,GeV spectral
  index $\GamLP$ (x-axis) against the log-parabola SED peak energy
  $\log\Epeak$ (y-axis), for the 2,630 labeled sources. BLL and FSRQ
  separate along a diagonal locus consistent with the classical blazar
  sequence \citep{Fossati1998}; PSR and MSP cluster separately owing
  to their intrinsically curved, cutoff-dominated spectra. This
  provides direct physical grounding for the $\log\Epeak$-driven
  improvement quantified in Fig.~\ref{fig:feat_compare}.}
\label{fig:blazarseq}
\end{figure}

\subsubsection{Robustness of the feature-set gain}
\label{sssec:robustness_results}

Because a single CV run or hold-out split can be sensitive to the
particular random partition, we repeated both evaluations following
the protocol of Section~\ref{ssec:robustness}. The repeated 10-fold CV
(5 seeds $\times$ 10 folds, 50 paired fold-level estimates) gives
$\BAcc = 84.7\pm3.6\,\%$ (7-feature) versus $86.1\pm3.2\,\%$
(10-feature), a gain of $+1.3$ percentage points with
$p = 8.8\times10^{-4}$ (paired Wilcoxon) -- a smaller point estimate of
the gain than the single-seed CV reported above, but a \emph{more}
significant one, because pooling over seeds reduces the fold-to-fold
variance that otherwise dominates the uncertainty. Independently, 15
repeated stratified 70/30 hold-out splits give test-set
$\BAcc = 85.4\pm2.1\,\%$ (10-feature) versus $84.3\pm2.0\,\%$
(7-feature), $p = 0.022$. Both resampling schemes therefore confirm
the direction and statistical significance of the original single-split
result while giving a more conservative, better-characterised estimate
of its magnitude ($+1.1$ to $+1.3$ percentage points, rather than the
$+2.3$ percentage points measured on the single original split, which
we attribute to a fold-partition that happened to be mildly favourable
to the 10-feature set). We report both the original and repeated
estimates for transparency; the repeated estimates should be regarded
as the more reliable characterisation of the true effect size.

\subsection{Classifier performance on the test set}
\label{sec:results_clf}

Table~\ref{tab:supervised} summarises the test-set results for all
three classifiers under the 10-feature set.

\begin{table}
\caption{Test-set performance ($N_{\rm test} = 789$) for three
  classifiers with the 10-feature set. $\BAcc$: balanced accuracy.
  AUC: macro one-vs-rest. CV\,$\BAcc$: 10-fold CV balanced accuracy
  (mean\,$\pm$\,std). Best values in bold.}
\label{tab:supervised}
\begin{tabular}{lcccc}
\hline
Classifier & Acc\,\% & $\BAcc$\,\% & AUC & CV\,$\BAcc$\,\% \\
\hline
RF  & \textbf{87.5} & \textbf{86.1} & \textbf{0.974} & $86.5\pm3.4$ \\
GB  & 88.2 & 86.4 & 0.969 & $86.3\pm2.1$ \\
MLP & 87.9 & 85.6 & 0.968 & $86.2\pm1.6$ \\
\hline
\end{tabular}
\end{table}

The RF achieves $\BAcc = 86.1\,\%$ and macro AUC \$= 0.974\$ on the
test set, outperforming GB and MLP modestly. The three classifiers
agree to within $1.7$ percentage points in $\BAcc$, indicating that
the performance is determined primarily by the feature set rather than
the choice of algorithm.

Fig.~\ref{fig:cm_roc} shows the confusion matrix and one-vs-rest ROC
curves for the RF. The dominant confusion channels are BLL$\to$MSP
($\sim$5\,\% of true BLL) and MSP$\to$BLL ($\sim$8\,\% of true MSP),
reflecting the genuine spectral overlap between low-synchrotron-peak
BLL and MSP at intermediate values of $\GamLP$ and $\sigs$. The
PSR class is the cleanest: its concentration at low $\GALL$ and high
$\sigs$ makes it almost entirely separable from the extragalactic
population. One-vs-rest AUC values are $\geq 0.96$ for all four
classes; the lowest is MSP (AUC\,$\approx 0.93$) due to the
BLL--MSP overlap.

\begin{figure}
\includegraphics[width=\columnwidth]{figs_final/fig5_cm_roc.pdf}
\caption{\emph{Left:} Normalised confusion matrix for the 10-feature
  RF on the held-out test set ($N = 789$). Row sums equal unity.
  \emph{Right:} One-vs-rest ROC curves for each class.}
\label{fig:cm_roc}
\end{figure}

\subsection{Feature importance}
\label{sec:results_fi}

Fig.~\ref{fig:fi} shows the mean decrease in Gini impurity for the
10-feature RF. The LP curvature parameter $\betaLP$ is the most
important single feature (18.5\,\%), followed by Galactic latitude
$\GALL$ (13.8\,\%), the spectral index at 1\,GeV $\GamLP$ (13.3\,\%),
the log flux uncertainty (11.7\,\%), and the LP peak energy
$\log\Epeak$ (11.3\,\%). The curvature significance $\sigs$, which
dominated the 7-feature set of \citet{Malyshev2023}, ranks lower
(8.1\,\%) in our 10-feature set because the addition of $\betaLP$,
$\GALL$, and $\log\Epeak$ absorbs part of its discriminative role.
The fractional variability $\Fvar$ contributes 9.1\,\%, and the
remaining features each contribute less than 6\,\%.

As a check against the known bias of Gini importance toward
continuous, high-cardinality features \citep{Breiman2001}, we also
compute permutation importance on the held-out test fold (15 repeats;
Section~\ref{ssec:robustness}). The two rankings agree on the top two
features -- $\betaLP$ and $\GamLP$ retain the largest importance under
both metrics -- but diverge further down the ranking: $\GALL$ and
$\log\Epeak$, ranked second and fifth by Gini importance, drop to
fourth and fifth by permutation importance, while $\sigs$ moves up
slightly. This reordering does not affect the paper's central claim
(the two leading features remain the LP curvature and the 1\,GeV
index under both metrics), but it does mean that the precise Gini-based
percentages in Fig.~\ref{fig:fi} should be read as indicative of
overall relevance rather than a precise ranking of the lower-importance
features, several of which carry overlapping information from
correlated flux and significance measurements.

The importance of $\betaLP$ is physically natural: this quantity
measures the rate at which the photon spectrum steepens with energy,
and is largest for sources with strong spectral curvature -- principally
FSRQs (curvature driven by the external Compton cooling break) and
pulsars (curvature driven by the magnetic-field exponential cutoff).
A large $\betaLP$ is therefore a near-necessary condition for both
pulsar-class sources and for FSRQ, making it the single most diagnostic
variable for this four-class problem.

\begin{figure}
\includegraphics[width=\columnwidth]{figs_final/fig6_feature_importance.pdf}
\caption{Random Forest feature importance (mean decrease in Gini
  impurity) for the 10-feature set, colour-coded by feature type.
  $\betaLP$ is the most important feature (18.5\,\%), followed by
  $\GALL$ and $\GamLP$. The three new features ($\GALL$, $\log\Epeak$,
  $\Fvar$) collectively account for 34.2\,\% of the total importance.}
\label{fig:fi}
\end{figure}

\subsection{Classification of unassociated and BCU sources}
\label{sec:results_unassoc}

We retrain the 10-feature RF on the full labeled sample of 2,630
sources and apply it to the 2,423 unassociated sources and 1,623 BCU
sources. To ensure the published catalogue is exactly reproducible,
the production model is trained with single-threaded tree construction
(\texttt{n\_jobs=1}); we verified bit-for-bit identical output across
independent reruns with the package versions listed in the Data
Availability section (this is in contrast to \textsc{scikit-learn}'s
default multi-threaded construction, \texttt{n\_jobs=-1}, which we
found is not guaranteed deterministic across environments even at
fixed \texttt{random\_state}, though it does not affect any of the
cross-validated or test-set statistics reported elsewhere in this
paper). For the unassociated sources, applying the OOD criterion
($\log\sqrt{TS} < 0.69$) flags 825 sources (34.0\,\%); their
classifications are retained but explicitly marked. Among the
remaining 1,598 in-distribution sources, 541 receive a high-confidence
classification ($P_{\rm max} \geq 0.80$): 274 BLL, 61 FSRQ, 74 PSR,
and 132 MSP.

Before adopting $P_{\rm max}\geq0.80$ as our high-confidence
criterion, we validated it directly: pooling out-of-fold predictions
across a 10-fold CV split on the labeled sample and binning by
predicted $P_{\rm max}$, the observed classification accuracy tracks
the predicted confidence closely across the full range
(Fig.~\ref{fig:calibration}), and sources with $P_{\rm max}\geq0.80$
are correctly classified 96.1\,\% of the time out-of-fold. The RF's
posterior probabilities are therefore reasonably well calibrated in
the high-confidence regime that matters for follow-up prioritisation,
which is a necessary (though not sufficient, given the covariate
shift documented in Section~\ref{sec:shift}) condition for the
threshold to be meaningful.

\begin{figure}
\includegraphics[width=\columnwidth]{figs_final/fig8_calibration.pdf}
\caption{Calibration of the 10-feature RF's posterior confidence,
  computed from pooled out-of-fold predictions across a 10-fold CV
  split on the labeled sample. The horizontal axis shows the mean
  predicted $P_{\rm max}$ within each of ten quantile bins; the
  vertical axis shows the observed fraction of correct classifications
  in that bin. The dashed diagonal marks perfect calibration; the
  dotted horizontal line marks the four-class chance level (0.25).
  The shaded histogram (right axis) shows the density of predicted
  $P_{\rm max}$ values, which is concentrated above 0.8. Axes are
  scaled to the actual range of the data (approximately 0.27--1.0)
  rather than the full [0,1] interval, since no predictions fall below
  the four-class floor.}
\label{fig:calibration}
\end{figure}

Table~\ref{tab:unassoc} summarises these numbers and Fig.~\ref{fig:sky}
shows the sky distribution of high-confidence candidates. The PSR
candidates are strongly concentrated toward the Galactic plane
($|b| < 15^\circ$), in good agreement with the known young pulsar
population. The BLL and FSRQ candidates are distributed
quasi-isotropically, as expected for an extragalactic AGN population.
No significant clustering or known over-density regions are visible in
either blazar class.

\begin{table}
\caption{Four-class classification of \fgl{} unassociated sources.
  HC: high-confidence ($P_{\rm max} \geq 0.80$, non-OOD). The HC
  sub-sample represents a conservative subset for follow-up
  prioritisation. Counts are from the deterministic production model
  (\texttt{n\_jobs=1}; see Data Availability).}
\label{tab:unassoc}
\begin{tabular}{lcccc}
\hline
 & \multicolumn{2}{c}{All unassoc.\ ($N=2423$)}
 & \multicolumn{2}{c}{HC non-OOD ($N=541$)} \\
Class & $N$ & \% & $N$ & \% \\
\hline
BLL  & 825 & 34.0 & 274 & 50.6 \\
FSRQ & 532 & 22.0 &  61 & 11.3 \\
PSR  & 395 & 16.3 &  74 & 13.7 \\
MSP  & 671 & 27.7 & 132 & 24.4 \\
\hline
OOD  & 825 & 34.0 & --- & ---  \\
\hline
\end{tabular}
\end{table}

\begin{figure}
\includegraphics[width=\columnwidth]{figs_final/fig7_sky_probabilities.pdf}
\caption{\emph{Left:} Sky positions of the 541 high-confidence
  ($P_{\rm max} \geq 0.80$, non-OOD) unassociated source candidates.
  PSR candidates (green) are concentrated near the Galactic plane.
  \emph{Right:} Distribution of $P_{\rm max}$ for high-confidence
  candidates in each class.}
\label{fig:sky}
\end{figure}

For BCU sources the classifier predicts 912 BLL (56.2\,\%), 596 FSRQ
(36.7\,\%), 31 PSR (1.9\,\%), and 84 MSP (5.2\,\%). The resulting
BLL:FSRQ ratio of 1.53:1 is somewhat higher than the 1.49:1 observed
among firmly identified blazars in \fgl{}, which is consistent with a
modest BLL excess in the optically faint BCU population where
spectroscopic follow-up has been less complete.

The full posterior probability table for all 2,423 unassociated
and 1,623 BCU sources, including Source\_Name, Galactic coordinates,
four class probabilities, $P_{\rm max}$, predicted class, and OOD
flag, is available as supplementary machine-readable data.

%% ── 6. Discussion ────────────────────────────────────────────────────
\section{Discussion}
\label{sec:disc}

\subsection{Where does the feature-set gain come from?}
\label{ssec:disc_feat}

The $+2.3$ percentage point improvement in balanced accuracy from
the 7-feature baseline to the 10-feature set warrants scrutiny, since
such gains can sometimes reflect data leakage or overfitting rather
than genuine discriminative information. Several considerations argue
against this concern.

First, the improvement is statistically significant on a paired
per-fold test (Wilcoxon $p = 0.029$) with ten independent folds,
and it is robust across three different classifiers (RF, GB, MLP all
improve by comparable amounts). Second, the gain is concentrated in
the FSRQ class, which is physically plausible: $\log\Epeak$ separates
BLL from FSRQ along the blazar sequence, and $\Fvar$ distinguishes
FSRQ high-amplitude variability. Third, none of the three new features
carries direct class-label information; in particular, $\log\Epeak$
is continuous and its median-imputation is class-blind.

The per-class F1 breakdown also provides a useful cross-check. The
PSR and MSP scores are essentially unchanged ($|\Delta F_1| < 0.005$),
which is expected since neither $\log\Epeak$ nor $\Fvar$ add
information to the spectral curvature / Galactic latitude signature
that already separates pulsars from blazars. The gain is cleanly
localised to the blazar sub-classification problem.

\subsection{The BLL--MSP confusion and its implications for follow-up}
\label{ssec:disc_confusion}

The dominant off-diagonal term in the confusion matrix is
BLL$\leftrightarrow$MSP: approximately 5\,\% of true BLL are
predicted MSP and vice versa in the held-out test. This confusion
arises from the genuine overlap of BLL and MSP in $\GamLP$ space
(Table~\ref{tab:class_stats}: BLL median $\GamLP = 1.84$, MSP median
$1.90$) combined with the limited training sample of MSP
($N = 179$, constrained by the completeness of the Third \fermi{}
Pulsar Catalog; \citealt{Smith2023}).

For practical follow-up scheduling, this BLL--MSP confusion implies
that candidates with moderate $\sigs$ ($< 4$), soft $\GamLP$
($1.7$--$2.0$), and low $\GALL$ ($< 20^\circ$) are the most
ambiguous. For these sources, dedicated radio observations (which
readily distinguish a pulsing MSP from a continuous BLL on timescales
of minutes) provide the most efficient resolution pathway.

\subsection{The elevated MSP fraction among all predictions}
\label{ssec:disc_msp}

The all-source MSP prediction fraction ($27.7\,\%$) substantially
exceeds the labeled MSP fraction ($6.8\,\%$). This inflation is
largely confined to low-significance OOD sources: the high-confidence
non-OOD MSP fraction ($24.4\,\%$) is still elevated relative to the
prior, but closer to the $\sim$20--25\,\% fraction expected if one
accounts for the known Galactic-plane excess in the unassociated
sample. Approximately 34\,\% of unassociated sources lie below the
training significance threshold -- the lowest-significance faint
sources for which the classifier has seen few training examples --
and these are predominantly routed to the MSP class, which has the
broadest spectral feature distributions among the four classes.

This behaviour is not a systematic error in the sense of a model
failure; it is an honest reflection of the irreducible uncertainty
associated with classifying sources that are genuinely less
characterised than the training set. The OOD flag is precisely
designed to signal this regime.

\subsection{Comparison with previous work}
\label{ssec:disc_compare}

Our test-set balanced accuracy of $87.3\,\%$ for the four-class
problem on \fgl{} improves on the $\sim$70\,\% of
\citet{Saz_Parkinson2016} on 3FGL and the $\sim$72\,\% of
\citet{Zhu2024} on 4FGL-DR3, primarily because of the additional
training data in \fgl{} and the improved feature set. For binary
BLL/FSRQ classification, our RF achieves $\sim$91\,\% accuracy,
consistent with the $91.2\,\%$ of \citet{Cooper2023} on 4FGL-DR3.

Our covariate shift analysis extends \citet{MalyshevRASTI2023} in two
ways: by computing Wasserstein distances (which are physically
interpretable and scale-aware) in addition to KS statistics, and by
quantifying the global shift via a binary shift classifier AUC.
The shift AUC of 0.985 is substantially higher than what would be
expected for a mild shift problem, confirming that importance
weighting or other distribution-matching techniques would be necessary
for any classifier that claims to recover calibrated class posterior
probabilities for the unassociated population -- a claim we
deliberately refrain from making.

\subsection{Limitations}
\label{ssec:limits}

Several limitations merit explicit acknowledgement.
(i)~Approximately 34\,\% of unassociated sources are OOD with respect
to the training significance distribution; their classifications carry
reduced reliability and should not be used for population-level
statistics without appropriate weighting.
(ii)~The MSP training sample is small ($N = 179$), limited by the
completeness of the current pulsar timing surveys. As
\citet{Smith2023} add further MSP identifications to future data
releases, classification performance for this class is expected to
improve markedly.
(iii)~We have not applied importance weighting to correct for
covariate shift in the posterior probabilities, because our shift
classifier AUC (0.985) is too high for the standard
Shimodaira--Huang--Gretton importance-weighting procedure to give
reliable density ratio estimates; at this level of shift, only
domain adaptation methods or direct retraining on counterpart-confirmed
sources would yield properly calibrated posteriors.
(iv)~The $\log\Epeak$ imputation for PL-fit sources (11.1\,\% of the
sample) uses the population median. While this is class-blind, it
injects the median $\Epeak$ into sources for which the LP curvature
is consistent with zero, potentially washing out a genuine absence of
curvature as a discriminative signal.

%% ── 7. Conclusions ───────────────────────────────────────────────────
\section{Conclusions}
\label{sec:conc}

We have extended the seven-feature four-class \fermi{}-LAT
classification framework of \citet{Malyshev2023} to a ten-feature set
by adding Galactic latitude, the LP SED peak energy, and the
fractional variability amplitude. On the full 2,630-source labeled
sample of \fgl{}, the extended set raises the 10-fold cross-validated
balanced accuracy from $84.5 \pm 3.3\,\%$ to $86.5 \pm 3.4\,\%$ on a
single representative split (Wilcoxon $p = 0.029$), and this gain is
confirmed, with tighter statistical significance but a more
conservative magnitude ($+1.1$ to $+1.3$ percentage points,
$p = 8.8\times10^{-4}$ over 50 repeated fold estimates), by repeated
cross-validation and independent hold-out resampling. The gain is
concentrated in the FSRQ class ($\Delta F_1 = +0.026$), is reproducible
across three independent classifiers (RF, GB, MLP), and is physically
interpretable in terms of blazar-sequence separation via $\log\Epeak$,
which we confirm directly with a $\GamLP$--$\log\Epeak$ projection.
Permutation importance corroborates the Gini-based feature ranking for
the two leading features ($\betaLP$ and $\GamLP$).

The covariate shift between labeled and unassociated sources is
severe: a binary shift classifier achieves AUC\,$= 0.985$, and all
ten features show significant distributional differences (KS
$p < 10^{-22}$). We flag 34.0\,\% of unassociated sources as
out-of-distribution based on detection significance and caution
against using their classifications for calibrated population
statistics. An out-of-fold calibration check confirms that our
adopted $P_{\rm max}\geq0.80$ high-confidence threshold corresponds to
an observed accuracy of 96.1\,\% among in-distribution sources.

Applying the 10-feature Random Forest to all 2,423 unassociated
\fgl{} sources, we identify 541 high-confidence, in-distribution
candidates (274 BLL, 61 FSRQ, 74 PSR, 132 MSP; $P_{\rm max} \geq 0.80$;
Table~\ref{tab:unassoc}) and 912/596 BLL/FSRQ-dominated candidates
among the 1,623 BCU sources. The production model is trained
deterministically (\texttt{n\_jobs=1}; see
Section~\ref{sec:results_unassoc} and Data Availability for the exact
package versions), so these counts are exactly reproducible. All
posterior probabilities, OOD flags, and predicted labels are provided
as a machine-readable table.

Future improvements to this framework could include: (a) kernel
mean matching or domain adaptation to correct for covariate shift
in the posterior probabilities; (b) incorporation of MSP candidates
from forthcoming pulsar timing surveys (FAST, MeerKAT) to expand
the training sample for the most uncertain class; and (c) extension
to a six-class problem separating high-synchrotron-peak from
low-synchrotron-peak BLL, for which the $\log\Epeak$ feature provides
a natural handle.

\section*{Acknowledgements}
D.M.\ acknowledges support from IISER Berhampur. P.M.\ acknowledges
support from the Department of Atomic Energy, Government of India
(Project RTI4002). This research made use of data provided by the
Fermi Science Support Center (FSSC). The \fgl{} catalogue was
accessed from \url{https://fermi.gsfc.nasa.gov/ssc/data/access/lat/14yr_catalog/}.
Software: \textsc{scikit-learn} \citep{Pedregosa2011},
\textsc{numpy} \citep{Harris2020}, \textsc{scipy} \citep{Virtanen2020},
\textsc{matplotlib} \citep{Hunter2007}, \textsc{pandas} \citep{McKinney2010}.
FITS-table parsing was implemented directly against the FITS binary
format rather than via \textsc{astropy.io.fits}, to keep the extraction
pipeline dependency-light.

\section*{Data availability}
The full probabilistic classification catalogue for all 2,423
unassociated and 1,623 BCU \fgl{} sources, together with all analysis
and validation code (including the robustness, calibration, and
permutation-importance scripts described in
Section~\ref{ssec:robustness}), a \texttt{requirements.txt} pinning
the exact package versions, and the trained model file, will be made
available in a public repository upon publication \citep{code2025}.
The \fgl{} catalogue is publicly available from the FSSC. All results were
produced with \textsc{python}~3.12, \textsc{scikit-learn}~1.8.0,
\textsc{numpy}~2.4.4, \textsc{scipy}~1.17.1, \textsc{pandas}~3.0.2,
and \textsc{matplotlib}~3.10.8. The production Random Forest
(Section~\ref{sec:results_unassoc}) is trained with single-threaded
tree construction (\texttt{n\_jobs=1}); we verified that this yields
bit-for-bit identical classification catalogues across independent
reruns with the pinned versions above, and recommend this setting to
anyone reproducing the published counts. All cross-validated and
test-set statistics reported elsewhere in this paper (balanced
accuracy, confusion matrix, Gini and permutation feature importances,
shift classifier AUC) are likewise exactly reproducible with the
pinned versions and were verified by independent reruns during
manuscript preparation.

%% ── Appendix ──────────────────────────────────────────────────────────
\appendix

\section{FITS binary-table reader}
\label{app:fits}

The standard Python FITS I/O library \textsc{astropy} was unavailable
in our compute environment. We therefore wrote a minimal pure-Python
FITS binary-table parser that reads IEEE 754 big-endian floats,
16- and 32-bit signed integers, and fixed-width ASCII strings directly
from the raw binary stream, interpreting the header cards
\texttt{NAXIS1}, \texttt{NAXIS2}, \texttt{TFIELDS}, \texttt{TTYPE$n$},
and \texttt{TFORM$n$} to compute per-column byte offsets and data types.
The implementation is validated against the published source counts and
spectral statistics of \fgl{} and is provided in the supplementary
code repository. Users with access to \textsc{astropy} are encouraged
to use that library instead.

%% ── References ────────────────────────────────────────────────────────
\bibliographystyle{mnras}
\begin{thebibliography}{99}

\bibitem[Abdo et~al.(2010)]{Abdo2010_3LAC}
Abdo A.~A., et al., 2010, ApJ, 715, 429

\bibitem[Abdo et~al.(2013)]{Abdo2013}
Abdo A.~A., et al., 2013, ApJS, 208, 17

\bibitem[Abdollahi et~al.(2020)]{Abdollahi2020}
Abdollahi S., et al., 2020, ApJS, 247, 33

\bibitem[Acero et~al.(2015)]{Acero2015_3FGL}
Acero F., et al., 2015, ApJS, 218, 23

\bibitem[Ackermann et~al.(2015)]{Ackermann2015_3LAC}
Ackermann M., et al., 2015, ApJ, 810, 14

\bibitem[Ajello et~al.(2020)]{Ajello2020}
Ajello M., et al., 2020, ApJ, 892, 105

\bibitem[Atwood et~al.(2009)]{Atwood2009}
Atwood W.~B., et al., 2009, ApJ, 697, 1071

\bibitem[Ackermann et~al.(2012)]{Ackermann2012}
Ackermann M., et al., 2012, ApJ, 753, 83

\bibitem[Ball et~al.(2023)]{Ballet2023}
Ballet J., Bruel P., Burnett T.~H., Lott B., 2023, arXiv:2307.12546

\bibitem[Bhat \& Malyshev(2022)]{Bhat2022}
Bhat A., Malyshev D., 2022, A\&A, 660, A87

\bibitem[Breiman(2001)]{Breiman2001}
Breiman L., 2001, Mach.\ Learn., 45, 5

\bibitem[Bhatta et~al.(2024)]{Bhatta2024}
Bhatta G., Gharat S., Borthakur A., Kumar A., 2024, MNRAS, 528, 976

\bibitem[Chiaro et~al.(2016)]{Chiaro2016}
Chiaro G., et al., 2016, MNRAS, 462, 3180

\bibitem[Coronado-Bl\'{a}zquez(2022)]{CoronadoBlazquez2022}
Coronado-Bl\'{a}zquez J., 2022, MNRAS, 515, 1807

\bibitem[Cooper et~al.(2023)]{Cooper2023}
Cooper N., Dainotti M.~G., Narendra A., Liodakis I., Bogdan M.,
  2023, MNRAS, 525, 1731

\bibitem[Doert \& Errando(2014)]{Doert2014}
Doert M., Errando M., 2014, ApJ, 782, 41

\bibitem[Finke et~al.(2021)]{Finke2021}
Finke T., Kr\"{a}mer M., Manconi S., 2021, MNRAS, 507, 4061

\bibitem[Fossati et~al.(1998)]{Fossati1998}
Fossati G., Maraschi L., Celotti A., Comastri A., Ghisellini G.,
  1998, MNRAS, 299, 433

\bibitem[Friedman(2001)]{Friedman2001}
Friedman J.~H., 2001, Ann. Statist., 29, 1189

\bibitem[Ghisellini \& Tavecchio(2008)]{GhiselliniTavecchio2008}
Ghisellini G., Tavecchio F., 2008, MNRAS, 387, 1669

\bibitem[Ghisellini et~al.(2017)]{Ghisellini2017}
Ghisellini G., Righi C., Costamante L., Tavecchio F.,
  2017, MNRAS, 469, 255

\bibitem[Guillemot et~al.(2016)]{Guillemot2016}
Guillemot L., et al., 2016, ApJ, 819, 11

\bibitem[Harris et~al.(2020)]{Harris2020}
Harris C.~R., et al., 2020, Nature, 585, 357

\bibitem[Hunter(2007)]{Hunter2007}
Hunter J.~D., 2007, Comput.\ Sci.\ Eng., 9, 90

\bibitem[Ho(1995)]{Ho1995}
Ho T.~K., 1995, in Proc. 3rd Int. Conf. Document Analysis and
  Recognition. p.~278

\bibitem[Kang et~al.(2019)]{Kang2019}
Kang S.-J., et al., 2019, ApJS, 245, 12

\bibitem[Kerby et~al.(2021)]{Kerby2021}
Kerby S., Kaur A., Falcone A.~D., Eskenasy R., Hancock F., Stroh M.~C.,
  Ferrara E.~C., Ray P.~S., Kennea J.~A., Grove E., 2021, ApJ, 923, 75

\bibitem[Kingma \& Ba(2015)]{Kingma2015}
Kingma D.~P., Ba J., 2015, in Int. Conf. Learning Representations
  (ICLR). arXiv:1412.6980

\bibitem[\protect\citeauthoryear{Kov\v{a}\v{c}evi\'{c} et al.}{2020}]{Kovacevic2020}
Kov\v{a}\v{c}evi\'{c} M., et al., 2020, MNRAS, 490, 4770

\bibitem[Lefaucheur \& Pita(2017)]{Lefaucheur2017}
Lefaucheur J., Pita S., 2017, A\&A, 602, A86

\bibitem[Luo et~al.(2020)]{Luo2020}
Luo S., Leung A.~P., Hui C.~Y., Li K.~L., 2020, MNRAS, 492, 5377

\bibitem[Malyshev \& Bhat(2023)]{Malyshev2023}
Malyshev D.~V., Bhat A., 2023, MNRAS, 521, 6195

\bibitem[Malyshev(2023)]{MalyshevRASTI2023}
Malyshev D.~V., 2023, RAS Techn.\ Instrum., 2, 735

\bibitem[Massey(1951)]{Massey1951}
Massey F.~J., 1951, J. Am. Stat. Assoc., 46, 68

\bibitem[McKinney(2010)]{McKinney2010}
McKinney W., 2010, in Proc.\ 9th Python in Science Conf., pp.~51--56

\bibitem[Mirabal et~al.(2012)]{Mirabal2012}
Mirabal N., et al., 2012, ApJ, 749, L20

\bibitem[Niculescu-Mizil \& Caruana(2005)]{NiculescuMizil2005}
Niculescu-Mizil A., Caruana R., 2005, in Proc.\ 22nd Int.\ Conf.\
  Machine Learning. pp 625--632

\bibitem[Nolan et~al.(2012)]{Nolan2012}
Nolan P.~L., et al., 2012, ApJS, 199, 31

\bibitem[Mukherjee \& Majumdar(2025)]{code2025}
Mukherjee D., Majumdar P., 2025, analysis and validation code repository
  (to be made public upon publication)

\bibitem[Pedregosa et~al.(2011)]{Pedregosa2011}
Pedregosa F., et al., 2011, J.\ Mach.\ Learn.\ Res., 12, 2825

\bibitem[Platt(1999)]{Platt1999}
Platt J., 1999, in Adv.\ Large Margin Classifiers. pp 61--74

\bibitem[Rumelhart et~al.(1986)]{Rumelhart1986}
Rumelhart D.~E., Hinton G.~E., Williams R.~J., 1986, Nature, 323, 533

\bibitem[Saz~Parkinson et~al.(2016)]{Saz_Parkinson2016}
Saz~Parkinson P.~M., et al., 2016, ApJ, 820, 8

\bibitem[Shimodaira(2000)]{Shimodaira2000}
Shimodaira H., 2000, J.\ Stat.\ Plan.\ Inference, 90, 227

\bibitem[Smith et~al.(2023)]{Smith2023}
Smith D.~A., et al., 2023, ApJS, 265, 56

\bibitem[Urry \& Padovani(1995)]{UrryPadovani1995}
Urry C.~M., Padovani P., 1995, PASP, 107, 803

\bibitem[Virtanen et~al.(2020)]{Virtanen2020}
Virtanen P., et al., 2020, Nature Methods, 17, 261

\bibitem[Wilcoxon(1945)]{Wilcoxon1945}
Wilcoxon F., 1945, Biometrics Bull., 1, 80

\bibitem[Zhu et~al.(2021)]{Zhu2021}
Zhu K.-R., Kang S.-J., Zheng Y.-G., 2021, Res. Astron. Astrophys., 21, 015

\bibitem[Zhu et~al.(2024)]{Zhu2024}
Zhu K.~R., Chen J.~M., Zheng Y.~G., Zhang L., 2024, MNRAS, 527, 1794

\end{thebibliography}

\end{document}
